\section{Conclusion}
\label{sec:conclusion}

This paper developed an optimization atlas in which algorithm names, computational mechanisms, evidence states, reproducibility, and research gaps are distinct but connected objects. The separation matters because optimization contains both genuine methodological diversity and structural overlap, while empirical conclusions remain sensitive to problem conditions, budgets, implementations, and parameter policies.

The mechanism-first representation organizes algorithms by the state they preserve, information they use, proposal geometry, and the memory, selection, adaptation, diversity, constraint, and model mechanisms that modify search. Within the declared scope, a 90-entry census provides the identity layer, while 42 source-backed fingerprints linked to 46 verified source records provide the representative mechanism layer. The counts serve different purposes and are not interpreted as 90 distinct mechanisms.

The evidence atlas separates documented, partial, contradictory, untested, and project-validated conditions. This prevents absence of study from being labelled as failure and prevents structural similarity from transferring empirical conclusions between algorithms. Theory is retained with its assumptions and guaranteed object rather than collapsed into a generic notion of convergence.

The project-generated BBOB study demonstrates the executable side of the framework. The frozen cross-dimensional campaign provides a validated condition-level characterization across four comparator mechanisms, 24 noiseless BBOB functions, dimensions 5--40, and normalized evaluation budgets. Algorithm-separated COCO logging, independent objective-call accounting, and held-out instances provide an auditable path from execution to interpretation. The experiment does not establish a universal ranking.

The evidence audit also constrains what the paper does not claim. Because the exact numerical decision rule required for a confirmatory binary failure frontier was not fully frozen before held-out execution, no post hoc exact frontier is reported. Because the resulting evidence does not establish a G4/G5 mechanism-level repair opportunity, no new optimizer or repair operator is introduced. Claim strength is limited by the strongest evidence specified and validated before interpretation.

The practical outcome is a decision sequence from condition to mechanism, evidence, resolving experiment, and claim. Future versions can add algorithms, benchmark domains, theoretical results, replications, or application evidence without changing the meaning of existing records. Unresolved cells can be converted into explicit evidence obligations whose minimal resolving experiment is known.

Optimization research does not require every unresolved condition to become a new algorithm. In many cases the more informative contribution is to determine what is known, what remains untested, which mechanism is implicated, and what evidence would distinguish a genuine limitation from an artifact of experimental design. The atlas provides a reproducible structure for making that distinction.
